\section{Evaluation}
\label{sec:evaluation}

\subsection{RQ1: Breadth, Agreement, and Explicit Failure Accounting}

Table~\ref{tab:coverage} reports cross-suite outcomes. \system agrees on 31 of 39
SuiteSparse matrices; all eight exceptions remain. All 17 PETSc cases complete,
seven OpenModelica models execute 535,728 embedded linear calls, and all 68 COPS
instances reach KKT comparison. Full-NLP coverage is narrower because resource
gates, terminal-only cases, and timeouts are retained.

\begin{table*}[t]
\centering
\caption{Cross-suite coverage from the authoritative overall report. ``Compared''
means a same-input agreement or performance comparison exists under that suite's
contract.}
\label{tab:coverage}
\begin{tabular}{lrrl}
\toprule
Suite & Attempted & Compared & Explicit exception \\
\midrule
SuiteSparse & 39 & 31 & 8 no common success \\
PETSc TS equations & 17 & 17 & 0 pending \\
OpenModelica MSL & 7 & 7 & 1 no linear subsystem \\
COPS KKT & 68 & 68 & 0 failed \\
COPS full NLP & 68 & 12 perf. & 46 terminal-only \\
PDEBench-derived & 7 & 7 & 0 unreported \\
\bottomrule
\end{tabular}
\end{table*}

These counts demonstrate integration breadth, not universal solver dominance. In
particular, a SuiteSparse failure and a COPS terminal-only solve remain failures of
the stronger performance claim even if the runtime rejects the unsuccessful path.

\subsection{RQ2: Verified PDE Workload Acceleration}

Figure~\ref{fig:pdebench} and Table~\ref{tab:pde-details} report 30 paired runs per
workload. Medians span $\PdeMinimumShort\times$--$\PdeMaximumShort\times$; every 95\%
lower bound exceeds one, and every workload wins at least \PdeMinimumRunWins{} runs.

\begin{table}[t]
\centering
\scriptsize
\caption{PDEBench-derived same-input comparisons. Speedups are medians over 30
paired process runs; brackets show fixed-seed bootstrap 95\% intervals. Each
\system time includes setup, candidate kernel, and original-equation gate.}
\label{tab:pde-details}
\begin{tabular}{@{}lrlr@{}}
\toprule
Workload & Solves & Baseline & Result [95\% CI] \\
\midrule
Advection & \PdeAdvectionSolves & upwind & $\PdeAdvectionMedian\times$ $[\PdeAdvectionLower,\PdeAdvectionUpper]$ \\
Darcy & \PdeDarcySolves & Jacobi--PCG & $\PdeDarcyMedian\times$ $[\PdeDarcyLower,\PdeDarcyUpper]$ \\
Burgers & \PdeBurgersSolves & scalar direct & $\PdeBurgersMedian\times$ $[\PdeBurgersLower,\PdeBurgersUpper]$ \\
Diffusion--Sorption & \PdeDiffusionSorptionSolves & tridiagonal & $\PdeDiffusionSorptionMedian\times$ $[\PdeDiffusionSorptionLower,\PdeDiffusionSorptionUpper]$ \\
Shallow Water & \PdeShallowWaterSolves & PCG & $\PdeShallowWaterMedian\times$ $[\PdeShallowWaterLower,\PdeShallowWaterUpper]$ \\
Navier--Stokes & \PdeNavierStokesSolves & 10-wkr PCG & $\PdeNavierStokesMedian\times$ $[\PdeNavierStokesLower,\PdeNavierStokesUpper]$ \\
CFD1D & \PdeCfdOneDSolves & PCG & $\PdeCfdOneDMedian\times$ $[\PdeCfdOneDLower,\PdeCfdOneDUpper]$ \\
\bottomrule
\end{tabular}
\end{table}

\begin{figure}[t]
\centering
\begin{tikzpicture}
\begin{axis}[
  ybar,
  width=\columnwidth,
  height=4.8cm,
  bar width=7pt,
  ymin=1,
  ymode=log,
  log basis y=10,
  ylabel={Verified speedup ($\times$)},
  xtick={0,1,2,3,4,5,6},
  xticklabels={Adv.,Darcy,Burg.,Diff.,SW,NS,CFD},
  x tick label style={rotate=38,anchor=east,font=\scriptsize},
  ymajorgrids,
  grid style={dotted},
]
\addplot+[
  fill=blue!55,
  draw=blue!75!black,
  error bars/.cd,
  y dir=both,
  y explicit,
]
  table[
    x expr=\coordindex,
    y=speedup,
    y error minus=error_minus,
    y error plus=error_plus
  ] {data/pdebench_repeated_timing.dat};
\end{axis}
\end{tikzpicture}

\caption{Verified end-to-end speedup over each workload's designated classical
baseline. The log axis exposes both the six low-single-digit improvements and the
single $\PdeMaximumShort\times$ result; error bars are bootstrap 95\% intervals over 30
run-level paired ratios.}
\label{fig:pdebench}
\end{figure}

Six workloads lie below $\PdeNonCfdMaximumShort\times$; only CFD1D exceeds
$100\times$, rejecting a universal two-order-of-magnitude claim. In the 30-pair
counterbalanced study, order ratios span \MinimumOrderRatio{}--\MaximumOrderRatio{}
(maximum shift \MaximumOrderShiftPercent\%).
\OrderIntervalsContainingOne{} of \OrderIntervalsTotal{} intervals include one;
Diffusion--Sorption alone shows a small nonzero effect, while all 420 runs remain
faster than baseline.

The mechanisms also differ. Advection and Burgers profit from equation-structured
batched recurrences and vectorized gates. Shallow Water and incompressible
Navier--Stokes use learned/structured transforms with independent residual checks.
The CFD workload has unusually favorable repeated structure and amortization. This
heterogeneity motivates routing by equation family and complete cost rather than a
single globally advertised expert.

\subsection{RQ3: Routing Versus Fixed and Hindsight Experts}

On the evaluated sparse family, the calibrated router selects SuperLU instead of a
fixed structured tridiagonal expert and achieves a paired median speedup of
$\CalibratedRouterSpeedup\times$ with a bootstrap 95\% interval of
$[\CalibratedRouterLower,\CalibratedRouterUpper]$ and a
\CalibratedRouterWinPercent\% paired
win rate. The maximum mixed QoI error is $1.75\times10^{-6}$, with no dangerous
misroutes in the report.

The gate-passing hindsight reference reveals remaining headroom. The calibrated plan's
median cost is $1.076\times$ the per-scenario hindsight oracle and lies within 5\%
of it on 41.3\% of samples. A leave-one-scenario-out nearest-neighbor selector has
only 53.1\% exact winner accuracy, yet reaches practical equivalence to the oracle
on substantially more requests. This difference confirms that exact label accuracy
is a poor primary metric when several gate-passing experts have near-equal complete cost.

The 5$\times$5 source-to-6$\times$6 held-out shift preserves the selected expert,
but not every timing relation. Among the eight experts that pass the original-
equation gate in both families, cost-rank Spearman correlation is
\ShiftGatePassingRankSpearman{} and maximum acceptance-probability calibration error is
\ShiftGatePassingMaxCalibration{}. The source-selected expert is
\ShiftSelectedRegret{}$\times$ the fastest gate-passing held-out expert by competition
median. If structurally ineligible experts are included before filtering, maximum
calibration error instead reaches \ShiftAllMaxCalibration{}. Thus structural
eligibility is a prerequisite for meaningful probability calibration, and preserved
winner identity does not imply zero complete-cost regret under size shift.

Without refitting, the conditioning/topology matrix retains a minimum paired-speedup
95\% lower bound of \ShiftMatrixMinimumLower{}, maximum filtered calibration error
\ShiftMatrixMaximumCalibration{}, and maximum selected-cost regret
\ShiftMatrixMaximumRegret{}$\times$. Topology changes
\TopologyGateStatusChanges{} expert gate statuses with zero mismatch or dangerous
misroute. Since the embedding remains identical, capability/gate filtering---not
embedding similarity alone---preserves selection; family and hardware transfer remain
untested.

In the held-out joint-policy study, quadratic and cubic profiles select budgets
$\JointShiftQuadraticBudget$ and $\JointShiftCubicBudget$. Across
$\JointShiftHeldoutScenarios$ disjoint scenarios, maximum oracle regret is
$\JointShiftMaximumRegret\times$, maximum pass-probability calibration difference is
$\JointShiftMaximumCalibration$, and all requests succeed with
$\JointShiftGateMismatches$ gate mismatches. The 32 first-stage rejections continue to
32 second-stage acceptances. These controlled families test profile transfer, not a
learned universal policy.

Across $\ConditionedActionCount$ actions and $\ConditionedHeldoutRequests$ held-out
requests, median cost error is $\ConditionedCostMedianError$; pass Brier/ECE are
$\ConditionedBrier$/$\ConditionedEce$. Conditioned regret is
$\ConditionedRegret\times$ versus $\StaticProfileRegret\times$ for static and
$\FixedActionRegret\times$ for fixed controls. The router forms
$\ConditionedDistinctPlans$ plans, changes from the mode on
$\ConditionedPlanChangePercent\%$ of requests, matches exhaustive enumeration, and
records $\ConditionedFallbacks$ terminal continuations and $\ConditionedGateMismatches$ gate
mismatches. This controlled result does not establish public-workload dominance.

Table~\ref{tab:suitesparse-final-routing} preserves both final SuiteSparse outcomes.
V5 uses immutable first-run observations; only its control-aware row is a
zero-execution counterfactual. V6 is the unique prefrozen first run.

\begin{table}[t]
\centering
\scriptsize
\caption{Final SuiteSparse routing outcomes. Lower regret is better. The v5
control-aware value is an offline replay over frozen observations (R); all other
entries are recorded first-run policies or controls.}
\label{tab:suitesparse-final-routing}
\begin{tabular}{lrr}
\toprule
Policy & V5 & V6 \\
\midrule
Current control-aware anchor & $\SuiteRouteVFiveReplayRegret^{\mathrm{R}}$ & $\SuiteRouteRegretPrecise$ \\
Raw conditioned model & $2.457077$ & $\SuiteRouteRawRegret$ \\
Training family-fixed & $4.489972$ & $\SuiteRouteRegretPrecise$ \\
Global fixed SuperLU & $\SuiteRouteVFiveFixedRegret$ & $1.001123$ \\
Static profile & $\SuiteRouteVFiveStaticRegret$ & $1.756465$ \\
\bottomrule
\end{tabular}
\end{table}

V5 falsifies a general repair claim: the current rule switches
$\SuiteRouteVFiveSwitchedRequests$/$\SuiteRouteVFiveRequests$ requests and remains
$\SuiteRouteVFiveReplayVsFixed\times$ global fixed. One family already shares the
global anchor; another harmful family action passes both cost and support tests.
This is within-support cost miscalibration, not gate failure.

V6 is favorable but narrow. Over $\SuiteRouteHeldoutMatrices$ matrices,
$\SuiteRouteHeldoutRequests$ requests, and five repetitions, regret is
$\SuiteRouteRegretPrecise\times$, $\SuiteRouteFixedImprovementPercent\%$ below fixed
and $\SuiteRouteStaticImprovementPercent\%$ below static. It ties family-fixed because
no specialization, adaptation, or transition is enabled. The guard removes the raw
model's $\SuiteRouteRawRegret\times$ excess; RHS-only and tolerance-only controls also
tie the anchor, so fine-grained request adaptation is not established.

Median/p95/max cost error is $\SuiteRouteCostMedianError$/
$\SuiteRouteCostPNinetyFiveError$/$\SuiteRouteCostMaxError$ and Brier/ECE is
$\SuiteRouteBrier$/$\SuiteRouteEce$, weak diagnostics despite near-oracle control.
All $\SuiteRouteSuccesses$ requests succeed with $\SuiteRouteFallbacks$ terminal
continuations and zero failure, gate, order, or DP mismatch. Interaction/order deltas
of $\SuiteRouteInteractionDelta$/$\SuiteRouteOrderDelta$ remain diagnostic; no
transition enters the frozen policy.

The nonlinear cascade supplies a counterexample to routing-as-free-speedup. On one
coupled family, the online Equation-MoE path is $0.975\times$ the classic baseline
with a confidence interval below one; on a cubic family it reaches $1.199\times$.
The router therefore needs workload qualification and negative outcomes, not only
aggregate win rates.

\subsection{RQ4: Learned Candidates, Correction, and Break-Even}

Table~\ref{tab:operators} compares the two held-out operator studies. The linear
study accepts all 6,400 requests with zero continuation or failure and achieves
$\LinearOperatorSpeedup\times$ $[\LinearOperatorLower,\LinearOperatorUpper]$ paired
online speedup. Its measured training cost breaks
even after \LinearOperatorBreakEven{} queries and remains amortized-positive at the
10,000-query projection.

The nonlinear operator's raw candidate is dramatically inaccurate: its maximum
normalized full-state and QoI errors are 597.5 and 469.3. The strict gate rejects
the raw candidate. Generic Newton correction then converts all 6,400 requests into
accepted same-accuracy solves and yields $\NonlinearOperatorSpeedup\times$ online
speedup with $[\NonlinearOperatorLower,\NonlinearOperatorUpper]$. A weighted
diagonal residual/Jacobi corrector remains at 0\% acceptance after 32 steps, showing
that correction is not a fungible post-processing layer.

The production-corrector frontier further separates the two families. The linear
candidate already reaches full acceptance at budget
\LinearMinimumCorrectionBudget{}, so extra correction is unnecessary on that family.
For the nonlinear family, budget zero accepts no request and budget one accepts only
\NonlinearBudgetOneAcceptancePercent\%; the minimum full-acceptance budget is
\NonlinearMinimumCorrectionBudget{}. At that operating point, candidate-inclusive
complete cost is \NonlinearBudgetTwoVsZeroCostRatio{}$\times$ the budget-zero path,
because avoiding the terminal classical solve more than pays for correction. All
sweep points complete without Runtime failure. These family-specific thresholds do
not imply one globally optimal correction budget.

The shared hybrid control makes this comparison end-to-end. On the linear family it
accepts \LinearSharedAccepted{} requests but reaches only
$\LinearSharedBaselineSpeedup\times$ versus the classical baseline; the verified
runtime is $\LinearOperatorVsSharedBaseline\times$
$[\LinearOperatorVsSharedBaselineLower,\LinearOperatorVsSharedBaselineUpper]$ faster.
On the nonlinear family, the same control rejects all candidates, executes
\NonlinearSharedFallbacks{} mandatory terminal continuations, and reaches only
$\NonlinearSharedBaselineSpeedup\times$; the verified runtime is
$\NonlinearOperatorVsSharedBaseline\times$
$[\NonlinearOperatorVsSharedBaselineLower,\NonlinearOperatorVsSharedBaselineUpper]$
faster. Both controls have zero failures and gate mismatches. Thus the gain is not
explained by merely adding any correction-and-continuation wrapper to the learned
candidate.

The native comparison executes official HINTS code, DeepONet, weights, and all
\HintsNativeCases{} public cases; both methods have zero failures and maximum common
residual \HintsNativeMaximumResidual{}. The default Router chooses
\texttt{structured-tridiagonal-direct-cpu-v1} rather than forcing a learned path and
is $\HintsNativeSpeedup\times$ $[\HintsNativeLower,\HintsNativeUpper]$ faster online;
setup-amortized speedup over 10,000 requests is $\HintsNativeAmortizedSpeedup\times$.
This supports complete-cost selection only for the official 1D Poisson workload, not
a claim that a \system{} latent operator beats HINTS or that the result transfers.

\begin{table}[t]
\centering
\footnotesize
\caption{Held-out operator studies. Speedup includes candidate construction,
correction, gate, and required runtime work.}
\label{tab:operators}
\begin{tabular}{lrr}
\toprule
Metric & Linear & Nonlinear \\
\midrule
Eval. requests & 6,400 & 6,400 \\
Accept/continue/failure & 6,400/0/0 & 6,400/0/0 \\
Validation cases & 64 & 64 \\
95\% error-rate upper bound & 4.57\% & 4.57\% \\
Online speedup & $\LinearOperatorSpeedup\times$ & $\NonlinearOperatorSpeedup\times$ \\
95\% interval & $[\LinearOperatorLower,\LinearOperatorUpper]$ & $[\NonlinearOperatorLower,\NonlinearOperatorUpper]$ \\
Shared-control speedup & $\LinearSharedBaselineSpeedup\times$ & $\NonlinearSharedBaselineSpeedup\times$ \\
Training wall time & \LinearOperatorTrainingMs{} ms & \NonlinearOperatorTrainingMs{} ms \\
Break-even queries & \LinearOperatorBreakEven{} & \NonlinearOperatorBreakEven{} \\
Max. state error & $1.61\times10^{-5}$ & $7.29\times10^{-5}$ \\
Raw candidate accepted & yes & no \\
\bottomrule
\end{tabular}
\end{table}

\subsection{RQ5: Intra-node Scaling and Heterogeneous Placement}

Figure~\ref{fig:ns-scaling} separates scaling of the original-equation gate from
the complete verified path. At ten workers, the gate reaches $2.068\times$ paired
speedup with a bootstrap 95\% interval of $[1.982,2.172]$, whereas the complete
path reaches only $1.207\times$ with $[1.100,1.354]$. Four workers reach
$1.179\times$ complete-path speedup with $[1.085,1.241]$; the two-worker interval
still includes one. Across all runs, the maximum original-equation relative
residual is $4.01\times10^{-16}$ and the maximum cross-solver relative error is
$1.37\times10^{-10}$.

\begin{figure}[t]
\centering
\begin{tikzpicture}
\begin{axis}[
  width=\columnwidth,
  height=0.52\columnwidth,
  xlabel={Worker count},
  ylabel={Paired speedup vs. one worker},
  xmin=0.5,
  xmax=10.5,
  ymin=0.85,
  xtick={1,2,4,8,10},
  grid=major,
  legend style={font=\scriptsize,at={(0.02,0.98)},anchor=north west},
  tick label style={font=\scriptsize},
  label style={font=\scriptsize},
]
\addplot+[
  mark=o,
  thick,
  error bars/.cd,
  y dir=both,
  y explicit,
] table[
  x=workers,
  y=total,
  y error minus=total_minus,
  y error plus=total_plus,
] {data/ns_parallel_scaling.dat};
\addlegendentry{Complete verified path}
\addplot+[
  mark=square*,
  thick,
  error bars/.cd,
  y dir=both,
  y explicit,
] table[
  x=workers,
  y=gate,
  y error minus=gate_minus,
  y error plus=gate_plus,
] {data/ns_parallel_scaling.dat};
\addlegendentry{Original-equation gate}
\end{axis}
\end{tikzpicture}

\caption{Intra-node worker scaling on 40 incompressible Navier--Stokes Helmholtz
solves over 30 repetitions. Error bars are fixed-seed bootstrap 95\% intervals of
paired speedup relative to one worker. Gate parallelism is useful, but the serial
setup and spectral kernel bound complete-path gain.}
\label{fig:ns-scaling}
\end{figure}

The gate result demonstrates a concrete parallel verification mechanism, while the
end-to-end result prevents it from being advertised as linear solver scaling.
Classical PCG reaches $3.69\times$ paired speedup at ten workers in the same study,
so the remaining \system bottleneck is the worker-independent spectral kernel and
setup rather than insufficient baseline parallelism. This is intra-node evidence
on one Apple M4 host, not a multi-node or cross-architecture result.

The gate-only study broadens the mechanism beyond the Navier--Stokes stencil.
Table~\ref{tab:gate-family-scaling} reports paired speedup over the same sequential
fused gate. At ten workers, the linear family reaches
$\LinearGateTen\times$ with
$[\LinearGateTenLower,\LinearGateTenUpper]$ and the nonlinear family reaches
$\NonlinearGateTen\times$ with
$[\NonlinearGateTenLower,\NonlinearGateTenUpper]$. Across 30 repetitions of 2,080
requests per family, decisions and
residuals exactly match the sequential gate. The different linear and nonlinear
scaling curves show that gate parallelism depends on arithmetic intensity and per-
request work; it is not a universal substitute for optimizing candidate generation
or correction.

\begin{table}[t]
\centering
\footnotesize
\caption{Gate-only worker scaling across two equation families. Values are paired
median speedups over one worker; all decision and residual mismatch counts are zero.}
\label{tab:gate-family-scaling}
\begin{tabular}{rrr}
\toprule
Workers & Linear gate & Nonlinear gate \\
\midrule
1 & $\LinearGateOne\times$ & $\NonlinearGateOne\times$ \\
2 & $\LinearGateTwo\times$ & $\NonlinearGateTwo\times$ \\
4 & $\LinearGateFour\times$ & $\NonlinearGateFour\times$ \\
8 & $\LinearGateEight\times$ & $\NonlinearGateEight\times$ \\
10 & $\LinearGateTen\times$ & $\NonlinearGateTen\times$ \\
\bottomrule
\end{tabular}
\end{table}

Table~\ref{tab:ns-batch} tests amortization at a fixed ten requested workers. From
12 to 160 real-data solves, complete-path per-solve time falls from $19.44\,\mu$s
to $7.89\,\mu$s. Normalized complete throughput reaches $2.432\times$ with a
bootstrap 95\% interval of $[2.256,2.521]$; gate throughput reaches
$2.427\times$ with $[2.168,2.725]$. The maximum original-equation relative
residual remains $5.71\times10^{-16}$ and the maximum cross-solver relative error
is $3.40\times10^{-10}$.

\begin{table}[t]
\centering
\footnotesize
\caption{Fixed-worker batch scaling, normalized to 12 solves. Intervals are from
30 paired repetitions; time is complete verified time per solve.}
\label{tab:ns-batch}
\begin{tabular}{rrl}
\toprule
Solves & Time/solve ($\mu$s) & Throughput scaling [95\% CI] \\
\midrule
12 & 19.44 & $1.000\times$ \\
20 & 13.70 & $1.337\times$ $[1.246,1.449]$ \\
40 & 10.71 & $1.741\times$ $[1.609,1.903]$ \\
80 & 8.59 & $2.224\times$ $[2.104,2.322]$ \\
160 & 7.89 & $2.432\times$ $[2.256,2.521]$ \\
\bottomrule
\end{tabular}
\end{table}

The batch curve confirms useful amortization but also saturates: a $13.3\times$
increase in request count yields only $2.43\times$ throughput scaling. Together
with Figure~\ref{fig:ns-scaling}, this identifies verification and setup/kernel
fractions as different resource-allocation targets rather than one generic
``parallel solver'' effect.

The generic device execution probe runs real Metal and Apple Neural Engine
candidates, verifies device/reference agreement, and accepts 64 GPU and two NPU
requests without continuation. However, PDEBench-specific probes show that device
backend support is not sufficient for performance. For Advection and Shallow Water,
Metal and ANE candidates take tens of milliseconds and fail the final $10^{-10}$
original-equation gate, while CPU FP64 structured paths complete in microseconds.
Burgers and Diffusion--Sorption device iterations similarly fail residual or
complete-cost criteria. The router therefore rejects those device paths.

The second linear operator replication is another important negative result. Its
paired bootstrap interval crosses one, so it does not establish stable acceleration
despite retaining same-accuracy behavior. This prevents a positive Poisson-family
result from being generalized to a different sparse topology.

Together, the negative paths validate the core design choice: heterogeneous
execution must be selected by verified end-to-end cost. A system that reports only
successful neural inference or forced accelerator kernels would reach the opposite
and incorrect conclusion on these workloads.
